\section{Exception and Interrupt Model}

VIV-32 uses a single exception and interrupt entry mechanism. Synchronous exceptions, software traps, system calls, and hardware interrupts all transfer control through the exception vector mechanism.

\begin{itemize}
\item Exceptions are identified by an exception cause code.
\item The exception cause code is written to \code{ecause}.
\item The exception return address or faulting instruction address is written to \code{epc}.
\item The associated fault address, when applicable, is written to \code{eaddr}.
\item Control transfers to a vector address derived from \code{evbase} and \code{ecause}.
\end{itemize}

\noindent
VIV-32 v0.1 defines machine mode only. User mode, supervisor mode, privilege transitions, nested privilege stacks, and memory protection are reserved for future versions.

\section{Exception Causes}

The \code{ecause} control register records the cause of the most recent exception or interrupt.

\begin{center}
\begin{tabular}{lll}
\hline
\code{ecause} value & Cause & Description \\
\hline
\code{0x00} & Reset & Reset entry \\
\code{0x01} & Illegal instruction & Undefined or reserved instruction encoding \\
\code{0x02} & Misaligned instruction fetch & Instruction fetch from a non-4-byte-aligned address \\
\code{0x03} & Misaligned data access & Misaligned halfword or word load/store \\
\code{0x04} & Software trap & \code{trap} instruction \\
\code{0x05} & System call & \code{syscall} instruction \\
\code{0x06} & Timer interrupt & Platform timer interrupt \\
\code{0x07} & External interrupt & Platform external interrupt \\
\code{0x08--0xFF} & Reserved & Reserved for future architectural or platform causes \\
\hline
\end{tabular}
\end{center}

\noindent
Cause values above \code{0xFF} are reserved in VIV-32 v0.1. Future versions may define additional architectural, platform, or implementation-specific exception causes.

\section{Exception Vector Calculation}

The \code{evbase} control register contains the base address of the exception vector table. Each exception cause has a 16-byte vector slot.

\begin{center}
\ensuremath{vector_address \assign evbase + (\code{ecause} \shleft 4)}
\end{center}

\noindent
Equivalently, the vector address is:

\begin{center}
\ensuremath{vector_address \assign evbase + \code{ecause} \times 16}
\end{center}

\noindent
The \code{evbase} register must contain a 16-byte-aligned address. Exception vector addresses are instruction addresses and must be 4-byte aligned.

\begin{itemize}
\item Each vector slot is 16 bytes.
\item Each vector slot can contain up to four 32-bit instructions.
\item A vector slot may contain a short handler directly.
\item A vector slot may contain a jump to a longer handler elsewhere.
\end{itemize}

\section{Exception Entry Sequence}

When an exception or interrupt is taken, the processor performs the following architectural actions:

\begin{enumerate}
\item Determine the exception cause.
\item Write the exception cause code to \code{ecause}.
\item Write the exception program counter value to \code{epc}.
\item Write the exception address to \code{eaddr}, if applicable.
\item Clear the interrupt-enable bit in \code{sr}.
\item Set the program counter to \ensuremath{evbase + (\code{ecause} \shleft 4)}.
\end{enumerate}

\noindent
Clearing the interrupt-enable bit prevents maskable interrupts from interrupting the initial exception handler sequence. Software may re-enable interrupts explicitly after saving any required state.

\section{Exception Program Counter}

The value written to \code{epc} depends on the exception or interrupt type.

\begin{center}
\begin{tabular}{lll}
\hline
Event type & \code{epc} value & Rationale \\
\hline
Illegal instruction & Faulting instruction address & Re-execution is possible after emulation or patching \\
Misaligned instruction fetch & Faulting instruction address & The fault is associated with the attempted fetch \\
Misaligned data access & Faulting instruction address & The faulting load/store may be emulated or reported \\
\code{trap} & Address of following instruction & Trap returns after the trap instruction \\
\code{syscall} & Address of following instruction & System call returns after the syscall instruction \\
Timer interrupt & Address of next instruction & Interrupt resumes normal execution \\
External interrupt & Address of next instruction & Interrupt resumes normal execution \\
\hline
\end{tabular}
\end{center}

\noindent
For synchronous faults, \code{epc} identifies the instruction that caused the exception. For software traps, system calls, and interrupts, \code{epc} identifies the instruction address at which execution should resume.

\section{Exception Address Register}

The \code{eaddr} control register records an address associated with an exception when such an address is meaningful.

\begin{center}
\begin{tabular}{lll}
\hline
Event type & \code{eaddr} value & Meaning \\
\hline
Illegal instruction & \code{0} & No separate fault address \\
Misaligned instruction fetch & Faulting fetch address & Misaligned instruction address \\
Misaligned data access & Effective address & Misaligned load/store address \\
\code{trap} & \code{0} & No separate fault address \\
\code{syscall} & \code{0} & No separate fault address \\
Timer interrupt & \code{0} & No separate fault address \\
External interrupt & Platform-defined or \code{0} & Optional platform interrupt source information \\
\hline
\end{tabular}
\end{center}

\noindent
If an exception or interrupt does not define a meaningful address, \code{eaddr} is written with zero unless otherwise specified by the platform.

\section{Status Register Interrupt State}

The \code{sr} control register contains the architectural condition flags and interrupt-control state.

\begin{center}
\begin{tabular}{lll}
\hline
Bit & Name & Meaning \\
\hline
\code{0} & \code{N} & Negative flag \\
\code{1} & \code{Z} & Zero flag \\
\code{2} & \code{C} & Carry flag \\
\code{3} & \code{V} & Overflow flag \\
\code{4} & \code{E} & Arithmetic error flag \\
\code{5} & \code{IE} & Interrupt enable bit \\
\code{6--31} & reserved & Reserved, must be zero in VIV-32 v0.1 \\
\hline
\end{tabular}
\end{center}

\noindent
When \code{IE} is set, maskable interrupts may be taken. When \code{IE} is clear, maskable interrupts are held pending or ignored according to platform interrupt-controller behaviour.

\begin{itemize}
\item The \code{ei} instruction sets \code{IE}.
\item The \code{di} instruction clears \code{IE}.
\item Exception and interrupt entry clears \code{IE}.
\item The \code{iret} instruction returns to the address stored in \code{epc}.
\end{itemize}

\noindent
VIV-32 v0.1 does not define an automatic previous-interrupt-enable bit. Software exception handlers are responsible for saving and restoring \code{sr} if precise restoration of interrupt state or flags is required.

\section{Interrupts}

Timer and external interrupts use the same entry mechanism as exceptions.

\begin{itemize}
\item Timer interrupts use exception cause \code{0x06}.
\item External interrupts use exception cause \code{0x07}.
\item Interrupts are maskable by the \code{IE} bit in \code{sr}.
\item Interrupts are taken only between instructions.
\item Interrupt entry writes the resume address to \code{epc}.
\item Interrupt entry clears \code{IE}.
\end{itemize}

\noindent
The exact source, priority, pending-bit behaviour, and acknowledgement mechanism for platform interrupts are defined by the platform specification.

\section{Software Trap}

The \code{trap} instruction raises a software trap exception.

\begin{itemize}
\item \code{ecause} is set to \code{0x04}.
\item \code{epc} is set to the address of the instruction following the \code{trap} instruction.
\item \code{eaddr} is set to zero.
\item The trap immediate is encoded in the faulting \code{trap} instruction and may be recovered by software from \ensuremath{epc - 4}.
\end{itemize}

\noindent
The interpretation of the trap immediate is defined by the execution environment. It may be used for debugging, monitors, firmware calls, or implementation-defined software services.

\section{System Call}

The \code{syscall} instruction raises a system-call exception.

\begin{itemize}
\item \code{ecause} is set to \code{0x05}.
\item \code{epc} is set to the address of the instruction following the \code{syscall} instruction.
\item \code{eaddr} is set to zero.
\end{itemize}

\noindent
The system-call ABI is defined separately by the VIV-32 ABI and execution-environment specifications.

\section{Illegal Instruction}

An illegal-instruction exception is raised when instruction decoding encounters an undefined, reserved, or otherwise invalid instruction encoding.

\begin{itemize}
\item \code{ecause} is set to \code{0x01}.
\item \code{epc} is set to the address of the illegal instruction.
\item \code{eaddr} is set to zero.
\end{itemize}

\noindent
If exception handling is unavailable, an illegal instruction halts the machine.

\section{Misaligned Access Exceptions}

A misaligned instruction fetch exception is raised when the processor attempts to fetch an instruction from an address that is not 4-byte aligned.

\begin{itemize}
\item \code{ecause} is set to \code{0x02}.
\item \code{epc} is set to the faulting instruction address.
\item \code{eaddr} is set to the misaligned fetch address.
\end{itemize}

\noindent
A misaligned data access exception is raised when a halfword or word load/store uses an effective address that does not satisfy the required alignment.

\begin{itemize}
\item \code{ecause} is set to \code{0x03}.
\item \code{epc} is set to the address of the faulting load/store instruction.
\item \code{eaddr} is set to the misaligned effective address.
\end{itemize}

\noindent
Byte loads and stores never raise misaligned data access exceptions.

\section{Exception Return}

Exception and interrupt handlers return using the \code{iret} instruction defined in the instruction set chapter.

\begin{itemize}
\item The program counter is loaded from \code{epc}.
\item Condition flags are not modified by \code{iret}.
\item \code{IE} is not automatically restored by \code{iret}.
\end{itemize}

\noindent
Because VIV-32 v0.1 does not define an automatic saved-status register, software handlers that need to restore the previous interrupt-enable state must save \code{sr} on entry and restore it before executing \code{iret}.

\section{Nested Exceptions and Interrupts}

Nested exceptions are permitted only when software explicitly enables them.

\begin{itemize}
\item Exception entry clears \code{IE}.
\item While \code{IE} is clear, maskable interrupts are not taken.
\item A handler may save \code{epc}, \code{ecause}, \code{eaddr}, and \code{sr}, then execute \code{ei} to permit nested interrupts.
\item If a nested exception occurs before software has saved the previous exception state, the architectural exception registers may be overwritten.
\end{itemize}

\noindent
Software handlers are responsible for preserving any general-purpose registers and control registers they need to restore.
